\chapter{Well-Founded Orders and Termination}\label{ch:wf}

Every previous chapter's rule was proved \emph{sound}: it does not change the value. This chapter
is about the other property a rewriting engine needs, that it \emph{stops}, and about the branch
of mathematics that proves it. The engine has no step budget; every evaluation terminates because
every rule makes a certain quantity smaller in a well-founded order, and the quantity was designed
rule by rule. The chapter relearns the theory --- well-founded relations, lexicographic products,
the connection to rewriting --- and then reads the engine's ordering as an application of it.

\section{Well-founded relations}

\begin{definition}
  A relation $\prec$ on a set $A$ is well-founded if every nonempty subset of $A$ has a
  $\prec$-minimal element: an $m \in S$ with no $s \in S$ satisfying $s \prec m$.
\end{definition}

\begin{theorem}
  $\prec$ is well-founded iff there is no infinite descending chain
  $a_0 \succ a_1 \succ a_2 \succ \cdots$.
\end{theorem}
\begin{proof}
  If there were such a chain, the set $\{a_i\}$ would have no minimal element. Conversely, if
  some nonempty $S$ has no minimal element, then every $s \in S$ has some $s' \in S$ with
  $s' \prec s$; choosing one for each $s$ (dependent choice) and starting anywhere gives an
  infinite descending chain.
\end{proof}

\begin{theorem}[Well-founded induction]
  If $\prec$ is well-founded on $A$ and $P$ is a property such that, for every $a$,
  ($P(b)$ for all $b \prec a$) implies $P(a)$, then $P(a)$ holds for every $a \in A$.
\end{theorem}
\begin{proof}
  If not, the set of counterexamples is nonempty and has a minimal element $a$; every $b \prec a$
  satisfies $P$, so $a$ does, a contradiction.
\end{proof}

$\mathbb{N}$ under $<$ is well-founded, and well-founded induction on it is strong induction.
That is also how Lean accepts a recursive definition: the recursive calls must be on arguments
smaller in some well-founded relation, and then the function is defined by well-founded
recursion. Every non-structural function in the engine --- the rewriter above all --- is
accepted this way.

\section{Lexicographic products}

\begin{theorem}
  If $\prec_1$ on $A$ and $\prec_2$ on $B$ are well-founded, then the lexicographic order on
  $A \times B$, in which $(a, b) \prec (a', b')$ iff $a \prec_1 a'$, or $a = a'$ and
  $b \prec_2 b'$, is well-founded.
\end{theorem}
\begin{proof}
  An infinite descending chain $(a_i, b_i)$ has non-increasing first components; since
  $\prec_1$ is well-founded, they are eventually constant; from that point the second components
  strictly descend, contradicting well-foundedness of $\prec_2$.
\end{proof}

By induction, any finite lexicographic product of well-founded orders is well-founded. This is the
theorem that lets a termination argument be built in \emph{tiers}: a first quantity that most
rules decrease, a second for the rules that leave the first unchanged, and so on. It is
\lc{Prod.lex} in Lean, and the engine's proof of well-foundedness is one application of it per
tier.

\section{Rewriting and reduction orders}

\begin{definition}
  A rewrite system terminates if there is no infinite sequence $t_0 \to t_1 \to \cdots$ of
  rewrite steps.
\end{definition}

\begin{theorem}
  A rewrite system terminates iff there is a well-founded order $\succ$ on terms with $t \succ t'$
  for every rewrite step $t \to t'$.
\end{theorem}
\begin{proof}
  If such an order exists, an infinite rewrite sequence would be an infinite descending chain.
  Conversely, if the system terminates, the transitive closure of $\to$ is itself such an order.
\end{proof}

The converse is useless in practice --- it takes termination as given --- so the work is always
to \emph{find} an order. Two families of orders handle most systems. A \emph{polynomial
interpretation} assigns each symbol a polynomial with natural coefficients and orders terms by
the value of the resulting expression; it is additive and easy to check, and it cannot handle a
rule that duplicates a subterm, since no polynomial in the size of $f$ pays for two copies of
$f$. A \emph{recursive path ordering} fixes a precedence on symbols and lets a symbol be replaced
by any term built from smaller symbols; it handles duplication and is the standard order for the
derivative rules, but it needs a single precedence, and Chapter~\ref{ch:order} shows that the
engine's simplifier needs $\mathtt{pow} > \mathtt{mul}$ for one rule and $\mathtt{mul} >
\mathtt{pow}$ for another.

\section{The simplifier's order}

The seven simplification rules of Part~I terminate under a weighted node count:

\begin{minted}{lean}
structure Weights where
  w : Expr → Nat
  pos : ∀ e, 1 ≤ w e
  head : ∀ e cs, w (withChildren e cs) = w e

def measure (W : Weights) : Expr → Nat
  | .num q => W.w (.num q)
  | .var x => W.w (.var x)
  | .add es => W.w (.add es) + measureList W es
  …

structure Rule (W : Weights) where
  name : String
  apply : Expr → Option RuleResult
  decreasing : ∀ e r, apply e = some r → measure W r.result < measure W e
\end{minted}

This is a polynomial interpretation with every symbol interpreted as a constant plus the sum of
its arguments. It is additive, which is what the rewriter's recursion needs (rewriting a child
by a smaller child makes the parent smaller), and each rule carries its decrease proof as a
field: the rule cannot be constructed without it. The weights --- numerals $2$, sums and products
$4$, powers $5$, variables and functions $8$ --- were chosen so that every rule decreases:
folding two numerals into one loses at least $2$; collecting $at + bt$ into $(a+b)t$ loses a
whole copy of $t$, which is where the guard that $t$ not be a bare numeral comes from.

\section{The pipeline's order}

The derivative rules duplicate. The product rule sends $\frac{d}{dx}(fg)$ to $f'g + fg'$: two
copies of $f$ and two of $g$ for one of each, and no additive measure survives. The engine's
answer is a lexicographic product of six natural-number tiers:

\begin{minted}{lean}
abbrev Mu := Nat × Nat × Nat × Nat × Nat × Nat
def μ (e : Expr) : Mu := (cmdCount e, d3Count e, litCount e, M e, size e, numCount e)

def MuLt (a b : Mu) : Prop :=
  Prod.Lex (· < ·) (Prod.Lex (· < ·) (Prod.Lex (· < ·) (Prod.Lex (· < ·) (Prod.Lex (· < ·) (· < ·))))) a b

theorem MuLt_wf : WellFounded MuLt :=
  (Prod.lex Nat.lt_wfRel (Prod.lex Nat.lt_wfRel (Prod.lex Nat.lt_wfRel (Prod.lex Nat.lt_wfRel
    (Prod.lex Nat.lt_wfRel Nat.lt_wfRel))))).wf
\end{minted}

The well-foundedness proof is the theorem of the second section, five times. The tiers are
counts of commands, of malformed derivatives, of nodes above matrix literals, then the weight
$M$, then size, then the magnitudes of integer numerals; each exists because a family of rules
could not be ordered by the tiers before it, and Chapter~\ref{ch:order} tells the story of each.
The tier that does the mathematical work is $M$:

\begin{minted}{lean}
def M : Expr → Nat
  | .num q => if q.isOne then 1 else if q.isInt then 2 else 4
  | .var _ => 10
  | .add es => ML es + (if es.isEmpty then 3 else 0)
  | .mul es => 2 + ML es + 2 * es.length
  | .pow b e => (M b + 1) * M e - 1
  | .fn "diff" [g, t] => 3 ^ (M g + 3) + M t
  | .fn _ es => 5 * ML es + 10
  | .matrix rows => 10 + MR rows
\end{minted}

This is a polynomial interpretation that is \emph{not} additive in two places, and each is a
theorem about a rule. A numeric exponent multiplies its base, so that $(b^m)^n \to b^{mn}$
decreases with symbolic $m$ while $x \cdot x \to x^2$ also decreases; and a derivative is
exponential in its body, so that the product rule --- which replaces one derivative of $fg$ by two
derivatives of the parts plus copies of the parts --- decreases, because $3^{M f + M g + \cdots}$
dominates $3^{M f + 3} + 3^{M g + 3} + 2(M f + M g)$ with room to spare. An exponential
interpretation of this kind is how a polynomial-style order can handle duplication after all: the
symbol being eliminated is made so heavy that any number of copies of its arguments are cheaper.

\begin{exbox}[$x \cdot x \to x^2$]
  $M(x \cdot x) = 2 + (10 + 10) + 2 \cdot 2 = 26$; $M(x^2) = (10 + 1) \cdot 2 - 1 = 21$. The
  per-factor charge of $2$ in the product is what pays. With $M(1) = 2$ instead of $1$, the rule
  $x \cdot x^a \to x^{a+1}$ stops decreasing; that is why $1$ weighs $1$.
\end{exbox}

\section{The innermost strategy}

The termination obligation for a pipeline rule is not that its output is smaller than its input.
It is conditional:

\begin{minted}{lean}
structure Ordered (rules : List PlainRule) : Prop where
  decreasing : ∀ r ∈ rules, ∀ e res, ChildrenNormal rules e →
    r.apply e = some res → res.error = none → MuLt (μ res.result) (μ e)
\end{minted}

The hypothesis \lc{ChildrenNormal} says no rule fires anywhere inside the children. It is needed
because the product rule duplicates its body, and if the body could still contain a command that
grows when it fires, no order of this shape survives --- the duplicated command grows twice. The
rewriter's \emph{innermost} strategy makes the hypothesis true at every firing: children are
normalized before the parent is tried. The proof of termination therefore carries a second
invariant through the recursion, ``the result is normal'', in the type of the rewriter's result.
This is the point where the theory of this chapter meets a strategy: the order is not on terms
but on terms whose children are already in normal form, and that is a legitimate well-founded
order because the strategy guarantees it is the only kind of term a rule ever sees.

\begin{minted}{lean}
theorem pipelineOrderedWith (norm : Norm) : Ordered (pipelineRulesWith norm)
\end{minted}

is the theorem, for every checker \lc{norm} the pipeline may be instantiated with, and its proof
is one decrease lemma per rule --- thirty-odd of them --- each an \lc{omega} goal after the
weights are unfolded. The termination line the notebook shows names it.

\section{What an order cannot afford}

An order is a budget of a different kind: it says which rewrites are affordable. Distribution
$a(b + c) \to ab + ac$ is not affordable together with $x \cdot x \to x^2$, and
Chapter~\ref{ch:cannot} shows that no interpretation of the engine's shape can make both
decrease; so expansion is a function (Chapter~\ref{ch:poly}), not a rule. $8^{1/2} \to 2 \cdot
2^{1/2}$ is not affordable at all, and the radicals of Chapter~\ref{ch:rad} are written as single
powers with a sixth tier to order them. And $\beta$-reduction is not affordable by any order,
because it does not terminate (Chapter~\ref{ch:lambda}); it is the one place with a budget. A
reader who has followed the book from the rationals to here has met every one of those decisions
as a piece of mathematics first and a design second, which is the order in which they were
relearned.

\begin{trybox}
\begin{minted}{text}
diff(x^3 * exp(x) * sin(x), x)
(a*b)^2 * (a*b)^3 * a
\end{minted}
  Open the capabilities panel and read the termination line. The first cell fires the product
  rule twice with nested chain and power rules, and there is no counter behind it.
\end{trybox}

\section{Exercises}

\begin{exercisebox}[Multiset order]
  Define the multiset extension of a well-founded order (a multiset decreases if one element is
  replaced by finitely many smaller ones) and prove it well-founded for multisets over $\mathbb{N}$.
  Which of the engine's rules would it order that the additive measure does not?
\end{exercisebox}

\begin{exercisebox}[Base 3]
  Compute $M$ on both sides of the chain rule $\frac{d}{dx}\sin u \to \cos u \cdot \frac{d}{dx} u$
  with the derivative's weight changed to $2^{M f + 3}$, find a $u$ for which it fails to
  decrease, and show that base $3$ works for every $u$.
\end{exercisebox}

\begin{exercisebox}[The obligation must be conditional]
  Give a term whose product-rule step \emph{increases} $\mu$ when the body contains an
  \texttt{expand} command, and explain why the innermost strategy never presents that term to
  the rule.
\end{exercisebox}

\begin{exercisebox}[Well-founded recursion in Lean]
  Write a Lean function on \lc{Nat × Nat} that recurses lexicographically and give its
  \lc{termination_by} clause. Then find, in \file{engine/MathEngine/Terminate.lean}, the
  \lc{termination_by} of the rewriter.
\end{exercisebox}
